\section{Research Gaps}
\label{sec:gaps}

This section critically analyzes major research gaps emerging from the synthesized literature, resolving key concerns regarding semantic ambiguity metrics, premature idea cementing in high-fidelity prototypes, Agile sprint integration, and threats to validity.

\subsection{Categorizing and resolving semantic ambiguity}
A recurring limitation of earlier empirical studies is their reliance on informal, subjective impressions of ``clarity'' or simple raw counts of elicited sentences. To establish scientific rigor, requirements ambiguity must be evaluated against formal linguistic classifications.

Drawing upon linguistic foundations in Requirements Engineering \cite{cascon2025ambiguity}, requirements ambiguity is categorized into four distinct classes: (1) \textit{lexical and semantic ambiguity}, where words or domain terms possess multiple meanings within the problem domain (e.g., conflicting interpretations of ``active user'' or ``transaction status''); (2) \textit{syntactic and structural ambiguity}, involving ambiguous grammatical structures such as dangling prepositions or unclear logical combinations of conditionals; (3) \textit{scope ambiguity}, arising when quantifiers (e.g., ``all'', ``each'', ``only'') leave operational boundaries indeterminate across system states; and (4) \textit{vagueness and incompleteness}, characterized by subjective, non-quantifiable adjectives (e.g., ``user-friendly'', ``appropriate response time'') that lack verifiable acceptance criteria.

Prototyping serves as an experiential filter across these linguistic categories: interactive visual mockups collapse pronoun and scope ambiguities by grounding interactions in concrete UI elements, while explicit screen layouts eliminate vagueness by forcing concrete inputs and measurable validation rules \cite{ferrari2016ambiguity, cascon2025ambiguity}. Furthermore, as emphasized by Spoletini et al. \cite{spoletini2018interview}, prototypes provide an objective reference anchor that reconciles divergent interpretations between analysts and stakeholders.

\subsection{Cost-benefit dynamics and high-fidelity risks}
While high-fidelity (Hi-Fi) interactive prototypes (e.g., clickable Figma screens, interactive front-ends) provide realistic user experience simulations, their premature adoption introduces significant organizational and cognitive risks.

A principal cognitive barrier is \textit{premature idea cementing}. When stakeholders and sponsors observe a polished, pixel-perfect interface early in the project, they anchor prematurely on the visual presentation \cite{bjarnason2023empirically, gupta2022strategic}. Consequently, stakeholders often become reluctant to propose structural alterations or challenge core business logic, mistakenly assuming the underlying architecture is finalized.

In parallel, high-fidelity prototypes induce \textit{premature release pressure}. Management may perceive surface visual completion as engineering readiness, exerting pressure on development teams to deliver prematurely or to reuse fragile prototype mockups in production environments.

From an operational cost perspective, low-fidelity paper prototypes require only a modest upfront investment (median preparation of 145 minutes \cite{rueda2020requirements}) and encourage stakeholders to critique and discard tentative concepts easily. Therefore, a prudent software engineering guideline is: \textit{deploy low-fidelity prototypes during exploratory elicitation and problem framing, transitioning to high-fidelity prototypes only during confirmatory validation before development}.

\subsection{Agile integration: mitigating late-stage rework}
A practical challenge in agile software development is that standard user stories (\textit{As a [Role], I want [Action], So that [Benefit]}) are often too concise, leaving critical UI layouts, data fields, and non-functional constraints unarticulated. This disconnect leads to substantial rework during sprint demonstrations.

To bridge this gap without compromising agility, recent empirical evidence supports coupling prototypes with user stories into \textit{Enriched User Stories} \cite{acm2022proposal}. In this workflow, the Product Owner and UI/UX designer work one sprint ahead of developers during backlog grooming. By attaching low-fidelity wireframes and visual acceptance criteria directly to user stories, teams clarify data fields and edge cases before coding begins, effectively reducing downstream sprint rework by 30\%--40\% \cite{acm2022proposal}.

\subsection{Threats to validity}
Following evidence-based software engineering guidelines, we evaluate four validity threats:

Regarding \textit{construct validity}, several studies treat gross requirement volume as a proxy for elicitation success. However, volume alone does not ensure completeness of non-visual backend logic; we mitigated this by synthesizing multi-dimensional metrics (completeness, correctness, typology). In terms of \textit{internal validity}, facilitator behavior and think-aloud overhead in interactive settings \cite{abad2018loud} may influence participant pacing and task execution.

For \textit{external validity}, primary controlled experiments \cite{rueda2020requirements, diaz2021gender, uv2023empirical} were conducted predominantly with student cohorts and academic domains; while offering high experimental control, findings should be corroborated across enterprise environments. Finally, regarding \textit{conclusion validity}, controlled experiments with modest sample sizes (e.g., 30 teams in Rueda et al. \cite{rueda2020requirements}) exhibit limited statistical power for detecting subtle completeness differences, requiring cautious interpretation of non-significant $p$-values.

\subsection{Identified research gaps}
Three primary gaps emerge from the synthesized literature. \textit{Gap 1 (Lack of Standardized Ambiguity Metric Applications):} Most comparative trials evaluate gross requirement volume rather than directly quantifying ambiguity reduction across linguistic categories (lexical, syntactic, scope, vagueness). \textit{Gap 2 (Absence of Empirical Longitudinal Rework Data):} While industrial frameworks indicate rework reductions from visual artifacts, longitudinal tracking of defect propagation across agile sprint deliveries remains scarce. \textit{Gap 3 (Academic vs. Industrial Contextual Divide):} Rigorous controlled experiments are primarily conducted in university cohorts, whereas industrial investigations rely heavily on qualitative interviews.
